
\subsection{弗洛尔上同调的挠与霍费尔距离：定理 J 的证明}
\label{subsec:6-5-3}

\noindent\textbf{由定理 6.1.25 推出定理 J。}
沿用定理~\ref{thm:6-1-25} 的记号。
该定理立即给出 $b'\ge b(\mu)$。
另一方面，\translatornote{原书在此直接由定理 6.1.25
推出 $b'\ge b(\mu)$；
但该定理只处理 $\lambda_{\downarrow i}>\mu$，
而 (6.1.23) 的 $b(\mu)$
计入 $\lambda_i=\mu$ 的临界项。
源证明未处理等号情形，
见疑误 A-0398。}
容易看出
\[
\#\bigl(\psi^{(0)}(L^{(0)})
\cap\psi^{(1)}(L^{(1)})\bigr)
\ge a'+2b'.
\]
结合 $a=a'$，即得定理 J。
\hfill$\square$

本小节继续考察弗洛尔上同调的挠部分
\liftIndex{I016}\liftIndex{I226}
与霍费尔距离\liftIndex{I218}。

\begin{example}{6.5.13}\label{ex:6-5-13}
设 $\Sigma$ 为黎曼曲面，
$L^{(0)},L^{(1)}$ 是其上的两个圆周。
假设恰有一个圆盘 $D^2\subset\Sigma$
对上边界算子分量
$\langle\delta[p],[q]\rangle$ 有贡献，
而不存在对
$\langle\delta[q],[p]\rangle$
有贡献的此类圆盘。
令 $\mu$ 为 $D^2$ 的面积。
容易看出
\[
HF\bigl((L^{(1)},0),(L^{(0)},0);
\Lambda_{0,\mathrm{nov}}\bigr)
\cong
\frac{\Lambda_{0,\mathrm{nov}}}
{T^\mu\Lambda_{0,\mathrm{nov}}}.
\]
容易找到哈密顿同痕 $\psi$，
使 $\psi(L^{(0)})\cap L^{(1)}=\varnothing$，
并且对任意 $\epsilon>0$，
它与恒等映射的霍费尔距离为
$\mu+\epsilon$。
因此，定理 J 在此例中是最优的
（见图~\ref{fig:6-5-1}）。
\end{example}

\begin{figure}[htbp]
\centering
\setlength{\unitlength}{0.74mm}
\begin{picture}(150,75)
\linethickness{0.35pt}
{\color{black!18}%
\polygon*(16,32)(18,30)(21,26)(25,23)(29,22)(33,23)
  (37,26)(40,29)(36,30)(32,31)(28,33)(23,34)(19,34)}
\qbezier(5,34)(4,75)(36,66)
\qbezier(36,66)(61,59)(75,56)
\qbezier(75,56)(94,81)(122,68)
\qbezier(122,68)(149,64)(148,35)
\qbezier(148,35)(145,5)(115,12)
\qbezier(115,12)(93,19)(78,35)
\qbezier(78,35)(64,7)(39,11)
\qbezier(39,11)(8,11)(5,34)
\put(45,48){\oval(23,13)}
\put(112,48){\oval(26,13)}
\qbezier(16,32)(28,12)(40,29)
\qbezier(40,29)(63,44)(101,56)
\qbezier(101,56)(128,67)(124,71)
\qbezier(20,13)(15,23)(16,32)
\qbezier(16,32)(21,38)(40,29)
\qbezier(40,29)(48,21)(71,33)
\qbezier(71,33)(100,51)(111,53)
\qbezier(111,53)(128,51)(127,46)
\put(18,45){\vector(1,-2){7}}
\put(13,49){$\scriptstyle\mu$}
\put(16,8){$\scriptstyle L^{(0)}$}
\put(79,38){$\scriptstyle L^{(1)}$}
\multiput(21,13)(2.5,3){6}{\circle*{0.55}}
\multiput(112,53)(1.5,3){7}{\circle*{0.55}}
\end{picture}
\par\smallskip
{\renewcommand{\thefigure}{6.5.1}
\refstepcounter{figure}
{\normalfont\bfseries 图~\thefigure}\label{fig:6-5-1}}
\end{figure}

\noindent\textbf{定理 6.1.25 的证明。}
为简化记号，
先只考虑横截的拉格朗日对子
$(L^{(0)},L^{(1)})$，
且两者均无障碍。
一般的干净相交情形，
其中 $L^{(i)}$ 弱无障碍，
或经 bulk 形变后弱无障碍，
可以同样处理；
但须谨慎选取相交流形
$L^{(0)}\cap L^{(1)}$
的奇异链复形中
可数生成的子复形。
这一修改留给有兴趣的读者。

令 $\{\psi_\rho^{(j)}\}_\rho$
（$j=0,1$）与 $\mu$
如式~\eqref{eq:6-1-24}，
并记 $L^{(j)\prime}=\psi^{(j)}(L^{(j)})$。
沿用第 5.3.2 小节的记号，
考虑弱滤过 $A_\infty$ 双模同态
\begin{equation}
\begin{aligned}
\varphi:\bigl(&C(L^{(1)},L^{(0)};
\Lambda_{\mathrm{nov}}),
\mathfrak n_{k_1,k_0}^{J_t,\mathfrak s}\bigr)\\
&\longrightarrow
\bigl(C(L^{(1)\prime},L^{(0)\prime};
\Lambda_{\mathrm{nov}}),
\mathfrak n_{k_1,k_0}^{J'_t,\mathfrak s'}\bigr).
\end{aligned}
\tag{6.5.14}\label{eq:6-5-14}
\end{equation}
它立于
$\mathfrak g^{1,\psi^{(1)},
\{J^{(1)}_\rho\}_\rho,
\mathfrak s_1^{\operatorname{top}(\rho)}}$
与
$\mathfrak g^{0,\psi^{(0)},
\{J^{(0)}_\rho\}_\rho,
\mathfrak s_0^{\operatorname{top}(\rho)}}$
之上，能量损失为 $\mu$
（见定义~\ref{def:5-2-1}
及式~\eqref{eq:5-3-15}）。
因为假设 $L^{(0)}$ 与 $L^{(1)}$
横截，此处以 $C(L^{(1)},L^{(0)};
\Lambda_{\mathrm{nov}})$ 简记
$C_2(L^{(1)},L^{(0)};
\Lambda_{\mathrm{nov}})$。
定义双模同态
$\phi(x)=\varphi(e^{\mathbf b_1}\otimes x
\otimes e^{\mathbf b_0})$。
于是
\begin{equation}
\phi\circ\delta_{\mathbf b_1,\mathbf b_0}
=\delta_{\psi_*^{(1)}\mathbf b_1,
\psi_*^{(0)}\mathbf b_0}\circ\phi.
\tag{6.5.15}\label{eq:6-5-15}
\end{equation}
这里
$\psi_*^{(j)}\mathbf b_j
=\mathfrak g_*^{j,\psi^{(j)},
\{J^{(j)}_\rho\}_\rho,
\mathfrak s_j^{\operatorname{top}(\rho)}}
(\mathbf b_j)
\in\mathcal M(\psi^{(j)}(L^{(j)}))$。
由于 $\varphi$ 的能量损失为 $\mu$，
$\phi$ 诱导同态
\begin{equation}
\begin{aligned}
\phi:
T^{\lambda+\mu}
C((L^{(1)},L^{(0)});\Lambda_{0,\mathrm{nov}})
\longrightarrow
T^\lambda
C((L^{(1)\prime},L^{(0)\prime});
\Lambda_{0,\mathrm{nov}}).
\end{aligned}
\tag{6.5.16}\label{eq:6-5-16}
\end{equation}

考虑逆向同痕
\[
\{\psi_\rho^{(j)}\}_\rho^{-1}
=\{\psi_{1-\rho}^{(j)}
\circ(\psi^{(j)})^{-1}\}_\rho,
\]
它从恒等映射通向
$(\psi^{(j)})^{-1}$。
注意
\[
\bigl\|\{\psi_\rho^{(j)}\}_\rho\bigr\|
=\bigl\|\{\psi_\rho^{(j)}\}_\rho^{-1}\bigr\|;
\]
这里的范数表示哈密顿同痕族的霍费尔长度
（见式~\eqref{eq:5-3-40}）。
因此，同理得到同态
\begin{equation}
\begin{aligned}
\phi'':T^{\mu+\lambda}
C((L^{(1)\prime},L^{(0)\prime});
\Lambda_{0,\mathrm{nov}})
\longrightarrow
T^\lambda
C((L^{(1)},L^{(0)});\Lambda_{0,\mathrm{nov}}).
\end{aligned}
\tag{6.5.17}\label{eq:6-5-17}
\end{equation}
它定义为
\begin{equation}
\phi''(x)=\varphi''
\bigl(e^{\psi_*^{(1)}\mathbf b_1}
\otimes x\otimes
e^{\psi_*^{(0)}\mathbf b_0}\bigr).
\tag{6.5.18}\label{eq:6-5-18}
\end{equation}
其中 $\varphi''$ 是与逆向同痕
$\{\psi_\rho^{(j)}\}_\rho^{-1}$
关联的弱滤过 $A_\infty$ 双模同态。
为缩短下面两个微分的指标，记
\[
\widetilde{\mathbf b}_j
:=(\psi^{(j)})^{-1}_*
\bigl(\psi_*^{(j)}\mathbf b_j\bigr)
\qquad(j=0,1).
\]
于是有恒等式
\begin{equation}
\delta_{\widetilde{\mathbf b}_1,
\widetilde{\mathbf b}_0}\circ\phi''
=\phi''\circ
\delta_{\psi_*^{(1)}\mathbf b_1,
\psi_*^{(0)}\mathbf b_0}.
\tag{6.5.19}\label{eq:6-5-19}
\end{equation}
由定理~\ref{thm:4-1-3}，
$\widetilde{\mathbf b}_j$
与 $\mathbf b_j$ 规范等价。
因此，命题~\ref{prop:5-2-37}
给出同态
\[
i:C((L^{(1)},L^{(0)});
\Lambda_{0,\mathrm{nov}})
\longrightarrow
C((L^{(1)},L^{(0)});
\Lambda_{0,\mathrm{nov}}),
\]
满足
\begin{equation}
i\circ\delta_{\widetilde{\mathbf b}_1,
\widetilde{\mathbf b}_0}
=\delta_{\mathbf b_1,\mathbf b_0}\circ i,
\tag{6.5.20}\label{eq:6-5-20}
\end{equation}
及
\begin{equation}
i\equiv\mathrm{id}
\pmod{\Lambda_{0,\mathrm{nov}}^+}.
\tag{6.5.21}\label{eq:6-5-21}
\end{equation}
置 $\phi'=i\circ\phi''$。
由式~\eqref{eq:6-5-19}、
\eqref{eq:6-5-20} 得
\begin{equation}
\phi'\circ
\delta_{\psi_*^{(1)}\mathbf b_1,
\psi_*^{(0)}\mathbf b_0}
=\delta_{\mathbf b_1,\mathbf b_0}
\circ\phi'.
\tag{6.5.22}\label{eq:6-5-22}
\end{equation}

\begin{lemma}{6.5.23}\label{lem:6-5-23}
对每个 $\lambda\ge0$，复合
\[
\begin{aligned}
\phi'\circ\phi:
T^{2\mu+\lambda}
C((L^{(1)},L^{(0)});\Lambda_{0,\mathrm{nov}})
\longrightarrow
T^\lambda
C((L^{(1)},L^{(0)});\Lambda_{0,\mathrm{nov}})
\end{aligned}
\]
与同样源靶的包含映射上链同伦。
\end{lemma}

\begin{proof}
对 $s\in[0,1]$、$T>0$，
取光滑函数 $\chi_{s,T}$，
满足以下条件：
\begin{align}
\chi_{s,T}
&\text{在 }(-\infty,0]\text{ 上非降，
在 }[0,\infty)\text{ 上非增，}
\tag{6.5.24.1}\label{eq:6-5-24-1}\\
\chi_{s,T}(t)&\equiv s
\quad(t\in[-T,T]),
\tag{6.5.24.2}\label{eq:6-5-24-2}\\
\chi_{s,T}(t)&\equiv0
\quad(|t|\text{ 充分大}).
\tag{6.5.24.3}\label{eq:6-5-24-3}
\end{align}
对每组 $s,T$，以
\begin{equation}
\psi_{\sigma,s,T}^{(j)}
=\psi_{\chi_{s,T}(\sigma)}^{(j)}
\tag{6.5.25}\label{eq:6-5-25}
\end{equation}
定义哈密顿形变
$\{\psi_{\sigma,s,T}^{(j)}\}_\sigma$。
注意
\begin{equation}
\bigl\|\{\psi_{\sigma,s,T}^{(j)}\}_\sigma\bigr\|
\le2\mu.
\tag{6.5.26}\label{eq:6-5-26}
\end{equation}

对每个 $s,T$，
沿参数 $\sigma$ 取这族同痕的
时间有序积（见第 5.3.2 小节），
得到弱滤过 $A_\infty$ 双模同态；
它诱导
\[
\phi_{s,T}:
C((L^{(1)},L^{(0)});\Lambda_{\mathrm{nov}})
\longrightarrow
C((L^{(1)},L^{(0)});\Lambda_{\mathrm{nov}}).
\]
由命题~\ref{prop:5-3-45}
及式~\eqref{eq:6-5-26} 得
\translatornote{原书实体页 400 的式 (6.5.27)
把源端写成 $T^{\mu+\lambda}C$；
但紧前的长度上界为 $2\mu$，
命题 5.3.45 只保证滤过下降不超过此长度，
故此处与引理陈述一致写作
$T^{2\mu+\lambda}C$，
见疑误 A-0396。}
\begin{equation}
\phi_{s,T}
\bigl(T^{2\mu+\lambda}
C((L^{(1)},L^{(0)});
\Lambda_{0,\mathrm{nov}})\bigr)
\subseteq
T^\lambda
C((L^{(1)},L^{(0)});
\Lambda_{0,\mathrm{nov}}).
\tag{6.5.27}\label{eq:6-5-27}
\end{equation}
令 $s=1$。
由标准的胶合论证
（见定理~\ref{thm:5-3-39} 的证明），
在复形的子模
\[
T^{2\mu}C((L^{(1)},L^{(0)});
\Lambda_{0,\mathrm{nov}})
\]
上有
\begin{equation}
\lim_{T\to\infty}\phi_{1,T}
=\phi''\circ\phi.
\tag{6.5.28}\label{eq:6-5-28}
\end{equation}
取足够大的 $T$，
考察双参数族
$\{\psi_{\sigma,s,T}^{(j)}\}_{s,\sigma}$。
注意 $\chi_{0,T}(t)\equiv0$。
沿 $\sigma$ 取时间有序积，
沿 $s$ 取逐时积，
得到 $\phi_{1,T}$ 与 $\phi_{0,T}$
之间的滤过 $A_\infty$ 双模同伦
（见第 5.3.3 小节）。
其中 $\phi_{0,T}$ 是恒等映射，
而 $\phi_{1,T}$ 收敛于
$\phi''\circ\phi$。
这里 $\phi$ 如式~\eqref{eq:6-5-14}
所设，$\phi''$ 如式~\eqref{eq:6-5-18}
所设。
结合 $\phi,\phi'$ 的定义、
式~\eqref{eq:6-5-26} 与
命题~\ref{prop:5-3-45}，
便得到引理的结论。
\end{proof}

\begin{remark}{6.5.29}\label{rem:6-5-29}
事实上，略微修改能量估计可证明：
复合 $\phi'\circ\phi$
还诱导同态
\[
\phi'\circ\phi:
T^{\mu+\lambda}
C((L^{(1)},L^{(0)});\Lambda_{0,\mathrm{nov}})
\longrightarrow
T^\lambda
C((L^{(1)},L^{(0)});\Lambda_{0,\mathrm{nov}}),
\]
它与包含映射上链同伦。
不过，下文不使用这一加强结论。
\end{remark}

取 $\lambda\ge2\mu$。
于是有上链映射
\begin{equation}
\begin{aligned}
\phi:\quad
&\bigl(T^\lambda C((L^{(1)},L^{(0)});
\Lambda_{0,\mathrm{nov}}),
\delta_{\mathbf b_1,\mathbf b_0}\bigr)\\
&\longrightarrow
\bigl(T^{\lambda-\mu}
C((L^{(1)\prime},L^{(0)\prime});
\Lambda_{0,\mathrm{nov}}),\\
&\hspace{6em}
\delta_{\psi_*^{(1)}\mathbf b_1,
\psi_*^{(0)}\mathbf b_0}\bigr),
\end{aligned}
\tag{6.5.30.1}\label{eq:6-5-30-1}
\end{equation}
以及
\begin{equation}
\begin{aligned}
\phi':\quad
&\bigl(T^{\lambda-\mu}
C((L^{(1)\prime},L^{(0)\prime});
\Lambda_{0,\mathrm{nov}}),\\
&\hspace{6em}
\delta_{\psi_*^{(1)}\mathbf b_1,
\psi_*^{(0)}\mathbf b_0}\bigr)\\
&\longrightarrow
\bigl(T^{\lambda-2\mu}
C((L^{(1)},L^{(0)});
\Lambda_{0,\mathrm{nov}}),
\delta_{\mathbf b_1,\mathbf b_0}\bigr).
\end{aligned}
\tag{6.5.30.2}\label{eq:6-5-30-2}
\end{equation}
它们的复合 $\phi'\circ\phi$
与包含映射
\[
\phi:
\bigl(T^\lambda C((L^{(1)},L^{(0)});
\Lambda_{0,\mathrm{nov}}),
\delta_{\mathbf b_1,\mathbf b_0}\bigr)
\longrightarrow
\bigl(T^{\lambda-2\mu}
C((L^{(1)},L^{(0)});
\Lambda_{0,\mathrm{nov}}),
\delta_{\mathbf b_1,\mathbf b_0}\bigr)
\]
上链同伦。
下面将用到一个代数引理。

\begin{lemma}{6.5.31}\label{lem:6-5-31}
设 $D$ 是有限生成的
$\Lambda_{0,\mathrm{nov}}^{(0)}$-模，
且
\[
D\cong
\bigl(\Lambda_{0,\mathrm{nov}}^{(0)}\bigr)^A
\oplus
\bigoplus_{i=1}^{B}
\frac{\Lambda_{0,\mathrm{nov}}^{(0)}}
{T^{c_i}\Lambda_{0,\mathrm{nov}}^{(0)}}.
\]
则有：
\begin{enumerate}[label=\textup{(\arabic*)},leftmargin=2.5em]
\item 对每个 $\lambda>0$，
$T^\lambda D$ 作为
$\Lambda_{0,\mathrm{nov}}^{(0)}$-模
所需的最少发生元数为
$A+\#\{i\mid c_i>\lambda\}$。
\item $T^\lambda D$ 的任意
有限生成子模 $D'$，\translatornote{原书实体页 400
在 (1) 把模类别误写为
$\Lambda_{0,\mathrm{nov}}$，
但 $D$ 只定义为
$\Lambda_{0,\mathrm{nov}}^{(0)}$-模；
在 (2) 又写“严格少于”所示数目，
与自由秩一的反例及紧后的证明相矛盾。
此处依型别写 $\Lambda_{0,\mathrm{nov}}^{(0)}$，
并将上界写为“至多”，
见疑误 A-0394、A-0397。}
至多由
$A+\#\{i\mid c_i>\lambda\}$
个元素生成。
\end{enumerate}
\end{lemma}

\begin{proof}
\textup{(1)} 来自
\[
T^\lambda D
\big/
T^\lambda\Lambda_{0,\mathrm{nov}}^{+(0)}D
\cong
\mathbb Q^{A+\#\{i\mid c_i>\lambda\}}.
\]
证明 \textup{(2)} 时，
对 $T^\lambda D$ 作同样论证即可，
所以只需考虑 $\lambda=0$。
置
$\widetilde D
=\bigl(\Lambda_{0,\mathrm{nov}}^{(0)}\bigr)^{A+B}$，
并取自然投影
$\pi:\widetilde D\to D$。
把 $D'$ 的发生元提升至
$\widetilde D$，
令 $\widetilde D'$ 为它们生成的子模。
则 $\widetilde D'$ 有限生成，
且 $\pi(\widetilde D')=D'$。
利用标准基
（引理~\ref{lem:6-3-2}），
可知 $\widetilde D'$ 至多由
$A+B$ 个元素生成；
故 $D'$ 也至多由
$A+B$ 个元素生成，得证。
\end{proof}

式~\eqref{eq:6-5-30-1}--
\eqref{eq:6-5-30-2} 诱导同态
\[
\begin{aligned}
\phi_*:
T^\lambda HF((L^{(1)},L^{(0)});
\Lambda_{0,\mathrm{nov}})
&\longrightarrow
T^{\lambda-\mu}
HF((L^{(1)\prime},L^{(0)\prime});
\Lambda_{0,\mathrm{nov}}),\\
\phi'_*:
T^{\lambda-\mu}
HF((L^{(1)\prime},L^{(0)\prime});
\Lambda_{0,\mathrm{nov}})
&\longrightarrow
T^{\lambda-2\mu}
HF((L^{(1)},L^{(0)});
\Lambda_{0,\mathrm{nov}}).
\end{aligned}
\]
其复合 $\phi'_*\circ\phi_*$
等于自然包含映射。
由于中间的
$T^{\lambda-\mu}
HF((L^{(1)\prime},L^{(0)\prime});
\Lambda_{0,\mathrm{nov}})$
由 $a+b'(\lambda-\mu)$ 个元素生成，
引理~\ref{lem:6-5-31} \textup{(2)}
表明，原上同调的
$T^\lambda$ 部分在
$T^{\lambda-2\mu}$ 部分中的像
至多由同样多的元素生成。
再用引理~\ref{lem:6-5-31} \textup{(1)}，
得到\translatornote{原书的 $b(\lambda)$
沿用 (6.1.23) 的非严格阈值
$\lambda_i\ge\lambda$，
但引理 6.5.31 的模计算要求
$c_i>\lambda$；
在临界等号处，
下面的式 (6.5.32)
尚需额外论证，见疑误 A-0398。}
\begin{equation}
a+b(\lambda)
\le a+b'(\lambda-\mu).
\tag{6.5.32}\label{eq:6-5-32}
\end{equation}

先在附加条件
$\lambda_{\downarrow i}>3\mu$
下证明定理~\ref{thm:6-1-25}。
按该定理，挠指数已按非增次序排列。
注意
“$\lambda_{\downarrow i}\ge\lambda$
且 $i\le b$”
等价于 $b(\lambda)\ge i$。
在式~\eqref{eq:6-5-32} 中取
$\lambda=\lambda_{\downarrow i}$，得
\[
b'\ge
b'(\lambda_{\downarrow i}-\mu)
\ge b(\lambda_{\downarrow i})
\ge i,
\]
故
$\lambda'_{\downarrow i}
\ge\lambda_{\downarrow i}-\mu$。
特别地，
$\lambda'_{\downarrow i}>2\mu$。
交换 $L^{(j)}$ 与 $L^{(j)\prime}$
并利用这个下界，
同样得到
$\lambda_{\downarrow i}
\ge\lambda'_{\downarrow i}-\mu$。
这就证明了
$\lambda_{\downarrow i}>3\mu$
时的定理~\ref{thm:6-1-25}。

为处理一般情形，
把哈密顿同痕分为足够短的若干段，
以便逐段应用上述论证。
仍用定理~\ref{thm:6-1-25} 的记号。
取分割
$\rho_0=0<\rho_1<\cdots<\rho_N=1$，
使
\begin{equation}
\lambda_{\downarrow i}
-\sum_{k=1}^{j}\mu_k
>3\mu_{j+1},
\qquad j=0,1,\ldots,N-1.
\tag{6.5.33}\label{eq:6-5-33}
\end{equation}
这里
\[
\begin{aligned}
\mu_i
={}&\operatorname{dist}
\bigl(\psi_{\rho_{i-1}}^{(1)},
\psi_{\rho_i}^{(1)}\bigr)\\
&+\operatorname{dist}
\bigl(\psi_{\rho_{i-1}}^{(0)},
\psi_{\rho_i}^{(0)}\bigr),
\end{aligned}
\]
而 $\operatorname{dist}$ 表示霍费尔距离。
由于端点距离不超过相应同痕段的长度，
$\sum_j\mu_j\le\mu<\lambda_{\downarrow i}$。
把两条同痕分得足够细，
使每段
$\mu_{j+1}<
(\lambda_{\downarrow i}-\mu)/3$，
便满足式~\eqref{eq:6-5-33}。
原书还写出以下精确赋值作为分段处方：
\[
\mu_{j+1}
\stackrel{\text{原书}}{=}
\min\left\{
\frac14\left(
\lambda_{\downarrow i}
-\sum_{k=1}^{j}\mu_k\right),
\mu-\sum_{k=1}^{j}\mu_k
\right\}.
\]
稍作扰动后，\translatornote{源书实体页 401
在将 $\mu_j$ 定义为端点霍费尔距离后，
仍写 $\sum_j\mu_j=\mu$ 并要求逐段
精确满足上式；
式 (6.1.24) 的 $\mu$ 却是给定路径的长度之和，
端点距离只保证总和不超过路径长度。
归纳仅需 (6.5.33)，
这里改用足够细的有限分段给出该条件，
见疑误 A-0395。}
可假设每个 $j=0,\ldots,N$ 时，
$\psi_{\rho_j}^{(0)}(L^{(0)})$
与 $\psi_{\rho_j}^{(1)}(L^{(1)})$
横截。

令 $\lambda_{\downarrow i}(j)$
表示
\[
HF\bigl(
\psi_{\rho_j}^{(1)}(L^{(1)}),
\psi_{\rho_j}^{(0)}(L^{(0)});
\Lambda_{0,\mathrm{nov}}
\bigr)
\]
的第 $i$ 个挠指数，
并按
$\lambda_{\downarrow i}(j)
\ge\lambda_{\downarrow i+1}(j)$
编号。
\translatornote{原书实体页 402
在此把两个拉格朗日子的次序写成
$(\psi^{(0)}(L^{(0)}),
\psi^{(1)}(L^{(1)}))$；
它与定理 6.1.25 及本节始终使用的
$(L^{(1)},L^{(0)})$ 次序相反。
此处依归纳起终点统一，
见疑误 A-0399。}
对 $j$ 归纳证明
\begin{equation}
\lambda_{\downarrow i}(j)
\ge\lambda_{\downarrow i}
-\sum_{k=1}^{j}\mu_k.
\tag{6.5.34}\label{eq:6-5-34}
\end{equation}
事实上，若对 $j$ 已成立，
式~\eqref{eq:6-5-33} 给出
$\lambda_{\downarrow i}(j)
>3\mu_{j+1}$。
于是可对第 $j+1$ 段应用
已在该附加条件下证明的
定理~\ref{thm:6-1-25}，
得到 $j+1$ 的式~\eqref{eq:6-5-34}。

因 $\lambda_{\downarrow i}(N)
=\lambda'_{\downarrow i}$，
可得
\[
\lambda'_{\downarrow i}
\ge\lambda_{\downarrow i}
-\sum_{k=1}^{N}\mu_k
\ge\lambda_{\downarrow i}-\mu.
\]
反向不等式的证明类似。
定理~\ref{thm:6-1-25} 证毕。
\hfill$\square$

把哈密顿同痕分段以控制能量损失的办法，
也曾用于 \textbf{[Ono95]}：
该文据此证明修正弗洛尔同调
是良定义的，
即不依赖滤过的选择。
